\documentclass[10pt,a4paper]{amsart}
\usepackage[margin=19mm]{geometry}
\usepackage{amsmath,amssymb,mathtools}
\usepackage[scr=rsfs,cal=euler]{mathalfa}
\usepackage{xfrac}
\usepackage[unicode,hidelinks,bookmarks=false]{hyperref}
\newcommand{\ie}{i.e.\ }
\newcommand{\cf}{cf.\ }
\newcommand{\resp}{respectively\ }
\newcommand{\Subsection}{\subsection}
\providecommand{\nobreakdash}{\nobreak\hbox{-}}
\setlength{\emergencystretch}{2em}
\multlinegap0pt
\usepackage{bm}
\usepackage{bbm}
\usepackage{dsfont}
\usepackage{xspace}
\usepackage{cancel}
\usepackage{mathtools}
\usepackage{amsthm}
\makeatletter
\@ifundefined{claim}{\newtheorem{claim}{Claim}}{}
\makeatother
\title[On a conjecture of Kolokolnikov on algebraic connectivity]{On a conjecture of Kolokolnikov on algebraic connectivity}
\author{Cheng Chi, Junjie Wang, Jiaxin Zheng}
\date{Source: arXiv:2608.09822v1 (CC BY 4.0)}
\begin{document}
\maketitle

\noindent\emph{Source and licence.} On a conjecture of Kolokolnikov on algebraic connectivity. Version arXiv:2608.09822v1 (CC BY 4.0), distributed under the Creative Commons Attribution 4.0 International licence, as recorded in the arXiv licence metadata for this exact version and shown on the abstract page https://arxiv.org/abs/2608.09822v1. Full rights notice: licenses/arxiv-2608.09822-CC-BY-4.0.txt.\par\smallskip

\noindent\textit{Editorial scope.} Counted unit: the Claim whose source label is \texttt{claim::1} in the article (source0.tex lines 275--277, source label \texttt{claim::1}), delivered together with the article's own complete original proof of it (source0.tex lines 279--334, one \texttt{proof} environment of 56 lines). No mathematics is written, repaired or re-derived: statement and proof are the article's own text line for line. Required context retained uncounted: none. Every definition, display and label of those lines is retained unchanged. Omitted from this package and not redistributed: 1--274, 278--278, 335--643. Nothing inside a retained excerpt is omitted or altered: every retained body is delivered exactly as the article's lines are written, so no omission receipt of any kind accompanies this package. Citations: the retained excerpts carry no citation command at all, so no bibliography entry of the article is transcribed and no reference is redistributed. Reference adaptation: none was needed and none was made; every reference inside the delivered unit resolves inside it, so no unresolved reference or citation is produced. Printed slips kept literal: no printed word, symbol, hypothesis or number of the counted statement or of its proof was altered. Layout and class: the article's own class is \texttt{article} with packages including booktabs, caption, mathrsfs, amsmath, amsfonts, amsthm, amssymb, bbding, graphics, multicol, graphicx, color, enumerate, rotating; the wrapper is the lane's pinned \texttt{amsart} editorial wrapper at 10pt on A4, which restates the theorem-like environments the retained text needs and adds the article's own macro definitions that the retained text needs (none), with extra wrapper packages (bm, bbm, dsfont, xspace, cancel, mathtools). The delivered numbering is the unit's own. Assets: no retained span contains a figure environment, a graphics-inclusion command, a PDF-page inclusion, a table or a TikZ picture; no image file is copied into this package, the package declares no asset dependency and no image file is redistributed with it. Licence: version arXiv:2608.09822v1 of this article is distributed under the Creative Commons Attribution 4.0 International licence, as recorded in the arXiv licence metadata for this exact version and shown on the abstract page https://arxiv.org/abs/2608.09822v1. A full rights notice is delivered at licenses/arxiv-2608.09822-CC-BY-4.0.txt. Characters: every retained excerpt is pure ASCII, so the delivered engine, XeLaTeX with its default Latin Modern text font, reports no missing character.
  \begin{claim}\label{claim::1}
    Let $A\subseteq V(G)$ be nonempty vertex set and $G[A]$ is connected. Assume that $\partial_G(A)\leq2|A|-\sigma$ and $\sigma\geq0$, then $6|A|\ge n+3+2\sigma$.
  \end{claim}

  \begin{proof}
    Let $B=N_G(A)$ and $C=V(G)\backslash N_G[A]$, where $|A|=a$ and $|B|=b$. Let $C_1,\cdots,C_k$ be the connected components of $G[C]$, and let $c_i=|C_i|$. Recall $\mu(C_i)=e(G[C_i])-c_i+1\geq 0$ be the cyclomatic number of the connected graph $G[C_i]$. There are no edges between $A$ and any $C_i$. Since $A$ is a good set, if some $C_i$ is also a good set, then there exists a good pair, which yields a contradiction. Hence, we have
    \begin{equation}\label{eq::5}
      \partial_G(C_i)\geq2c_i+1
    \end{equation}
    for $1\leq i\leq k$.

    For a vertex set $U$, define $\chi(U)=\sum_{v\in U}(d_G(v)-4)$. Since $e(G)=2n-4$, we have
    \begin{equation}\label{eq::6}
      \chi(V(G))=\sum_{v\in V(G)}(d_G(v)-4)=2e(G)-4n=-8.
    \end{equation}
    By inequality \eqref{eq::5},
    \begin{equation*}
      \begin{aligned}
        \chi(C_i)
         & =\sum_{v\in V(C_i)}(d_G(v)-4)           \\
         & =2e(G[C_i])+\partial_G(C_i)-4c_i        \\
         & =2(c_i-1+\mu(C_i))+\partial_G(C_i)-4c_i \\
         & =\partial_G(C_i)-2c_i-2+2\mu(C_i)       \\
         & \geq-1.
      \end{aligned}
    \end{equation*}
    Substituting these inequalities into \eqref{eq::6} gives
    $
      -8= \chi(V(G))=\chi(A)+\chi(B)+\sum_{i=1}^k\chi(C_i)
      \geq \chi(A)+\chi(B)-k,
    $
    and therefore
    \begin{equation}\label{eq::7}
      k\ge\chi(A)+\chi(B)+8.
    \end{equation}

    Using inequality \eqref{eq::5}, we have
    \begin{equation}\label{eq::8}
      e_G(B,C)=\sum_{i=1}^k\partial_G(C_i)\geq2|C|+k.
    \end{equation}
    On the other hand, summing degrees over $B$, we obtain
    $4b+\chi(B)=\sum_{v\in B}d_G(v)=2e_G(B)+\partial_G(A)+e_G(B,C)$.
    Consequently,
    \begin{equation}\label{eq::9}
      e_G(B,C)\leq 4b+\chi(B)-\partial_G(A).
    \end{equation}
    Combining \eqref{eq::7}--\eqref{eq::9} and using $|C|=n-a-b$ gives
    \begin{equation}\label{eq::10}
      6b\geq 2n-2a+\partial_G(A)+\chi(A)+8.
    \end{equation}
    Every vertex of $B$ has at least one neighbor in $A$, so $b\le \partial_G(A)$. Hence \eqref{eq::10} implies
    $
      6\partial_G(A)\geq 2n-2a+\partial_G(A)+\chi(A)+8$.
    Since $\chi(A)=2e_G(A)+\partial_G(A)-4a$,
    we obtain
    \begin{equation}\label{eq::11}
      4\partial_G(A)+6a\geq 2n+2e_G(A)+8.
    \end{equation}
    The graph $G[A]$ is connected, so $e_G(A)\geq a-1$. Substituting this into \eqref{eq::11} gives $4\partial_G(A)+4a\geq 2n+6$. Since $\partial_G(A)\leq2a-\sigma$, and therefore $4(2a-\sigma)+4a\ge2n+6$. Rearranging proves this claim.
  \end{proof}

\end{document}
